\documentclass[../main]{subfiles}
\begin{document}
\chapter*{Acuerdo Áulico}


\begin{table}[h!]
    \centering
    \begin{tabular}{|c|c|} \hline 
Materia: & Mecánica de los fluidos \\ \hline
Curso: &$5^{to} 3^{ra}$ \\ \hline
Docente: & Ferreyra Gustavo David\\ \hline
Horario: &Martes 9:15-11:25 \\ \hline
Área & Electromecánica \\ \hline
    \end{tabular}
\end{table}
 

\section{Introducción} 

El presente acuerdo tiene normas que tienen por objetivo organizar la presentación de trabajos prácticos, a fin de que el alumno sea consciente de la metodología con la que serán evaluados sus trabajos. Además, se busca, que el padre o tutor, esté al tanto de la evolución del alumno dentro del curso. 

\section{Consideraciones generales.} 

Nuestra Institución se ha caracterizado por poseer un nivel de exigencia alto y estándares de disciplina bastante estrictos, que han permitido que los egresados sean profesionales de conducta intachable y de amplias capacidades y conocimientos. Dichas virtudes les han abierto las puertas a trabajos y experiencias únicos.  

En el devenir de los tiempos modernos, es que debemos formalizar el trabajo y dejar acuerdos que permitan un trabajo docente transparente y claro para cada uno de los partícipes de la relación educativa. 

Por lo tanto, el presente acuerdo aúlico se enmarca en el Acuerdo Institucional de Normas de disciplina, de las resoluciones de DGE, y de todas las normativas vigentes. 

Los contenidos de toda la materia se encuentran debidamente avalados por el diseño Curricular Provincial.  

Se deja expresamente aclarado que la relación Docente-Alumno, estará basada en el respeto mutuo y en la plena convicción de que el aprendizaje se basa en la construcción permanente del conocimiento mediante técnicas y herramientas de investigación y trabajo colaborativo, que resignifiquen la función docente y el rol del alumno. En definitiva, se busca generar alumnos con espíritu crítico, capaces de formar equipos de trabajo productivos y que posean rasgos de liderazgo y de cooperación. 

 

\section{Acuerdos sobre el cursado.} 
 Sobre la conducta y trabajo en clases:
 El alumno deberá estar en el aula una vez que haya tocado el timbre. Si por algún motivo el docente se retrasa unos minutos, los alumnos deberán permanecer en el interior del aula guardando la conducta correspondiente. Si el alumno necesita un elemento de biblioteca o preceptoría, lo deberá buscar en el recreo. 

 Es también necesario el cuidado de la limpieza en el entorno, por la tanto, se cuidará no tirar papeles o restos de lápices o cualquier otro residuo que ensucie el piso del aula. “Lo que hagamos en el aula es un espejo de nuestras costumbres” 

En todo momento se cuidará el orden y la tranquilidad. Promoviendo el trabajo y los buenos tratos entre compañeros 

 

\section{Acuerdos sobre la Carpeta de Trabajos Prácticos (CTP). }

La carpeta de trabajos prácticos es una herramienta de trabajo y un documento de la trayectoria realizada por el alumno. Es necesario que el alumno sea el responsable del armado y confección de la misma, con los mayores criterios de responsabilidad y cumplimiento.  

El formato de la misma, se especifica en el presente documento. La justificación a las exigencias en cuanto a la CTP, están basadas en la unificación de criterios y formatos, lo que facilitará la evaluación de la misma.  

La presentación parcial de la CTP es obligatoria tantas veces como el docente la requiera, las presentaciones cuatrimestrales, serán con previo aviso, fijando una fecha de común acuerdo. 

 

\subsection{- Formato de la Carpeta de Trabajos Prácticos }

\begin{itemize}
    
\item Tamaño: A4 Normalizada 

\item Nombre de la materia 

\item Apellido y nombre del alumno 

\item Trabajo Práctico Nº  

\item Hoja Nº 

\end{itemize}
La carpeta deberá ser confeccionada a mano y en casos excepcionales y en acuerdo con el docente, con procesador de texto. En el desarrollo de la carpeta NO se exigirá letra técnica. Sin embargo, la realización de esta implica una buena presentación y permite una lectura más precisa y detallada, por lo tanto, una mejor nota. La estructura de la carpeta será: 
\begin{enumerate}
    
\item La primera hoja debe ser la portada, con todos los datos completos. 

\item La segunda hoja debe ser las presentes normas y deberá estar firmada por el alumno y padre o tutor. Dicha firma implica el conocimiento total de las presentes normas. 

\item La tercera hoja debe corresponder al programa. 

\item La cuarta hoja debe corresponder al Índice.  

\end{enumerate}
En lo relacionado con el formato general, todos los problemas que lo requieran deben llevar su gráfico o dibujo, el cual debe ser confeccionado prolijamente (usando regla; compás, etc.) y en tinta. El subrayado, recuadros, líneas de fracción, etc., deben ser hechos con regla. 

Al final de la carpeta se colocarán las evaluaciones escritas, las cuales deben estar firmadas por el padre o tutor al lado de la nota correspondiente. Las mismas deberán estar en la carpeta cada vez que se la requiera. 

La carpeta de Trabajos Prácticos debe ser correctamente presentada, es decir con sus hojas adecuadamente tomadas y no sueltas. Durante la presentación, la carpeta, se presentará en forma diferenciada del resto de las materias. La carpeta será parte de la nota final y servirá para lograr la aprobación o no del período correspondiente. 

La carpeta debe estar completa cada vez que el docente lo crea necesario. Los trabajos prácticos de investigación para ser presentados deben estar confeccionados en computadora según el siguiente formato: 
\begin{itemize}
    
\item Portada: datos de los alumnos participantes, docente, materia, tema, curso, año. 

\item Índice (en caso de ser necesario) 

\item Contenido – Desarrollo. (con rótulo correspondiente) 

\item Conclusiones. 

\item Bibliografía – Webs consultadas. 

\end{itemize}
 

\section{Acuerdos sobre el uso de tecnología en el aula. } 

Durante la clase, está terminantemente prohibido el uso de celulares y cualquier como elemento que fomente la distracción (fotos, videos, juegos, chats, etc). El no acatamiento de esta norma, hará pasible al alumno de ser sancionado disciplinalmente. En caso de requerir con urgencia el uso de un celular, se deberá pedir permiso al docente. El uso de las netbooks, notebooks o celulares queda definido y emplazado a servir de instrumento del aprendizaje y una herramienta útil. Queda totalmente prohibida la utilización de dichos equipos para recreación, juegos, chat o cualquier otra actividad que no construya el saber y que propicie la distracción del alumno. Queda a criterio del docente la sanción disciplinaria correspondiente a tal falta. 

Uno de los aspectos que se busca fortalecer es el trabajo con herramientas virtuales y de trabajo colaborativo. Para ello, se vuelve imprescindible trabajar responsablemente con todas las herramientas tecnológicas. 

 

En definitiva, los celulares, netbooks y notebooks únicamente pueden ser utilizados para ser parte del trayecto pedagógico y de la construcción del conocimiento. Para ello, se utilizarán softwares de diseño, aplicaciones, Google Classroom y simuladores, que nos permitirán mejorar notablemente la experiencia de aprendizaje en aula y en forma virtual. 

 

\section{Acuerdos sobre visitas o salidas educativas.} 

Las salidas o visitas son otra instancia más en el fortalecimiento del aprendizaje de cada alumno. Son recursos que propician el conocimiento y que permiten un acercamiento a la realidad laboral actual, y son posibilidades únicas que la condición de alumno otorga.  

Es por ello, que debemos participar de estos eventos con la mayor responsabilidad posible y apuntando a alcanzar el mayor nivel de adquisición de conocimientos permitido. 

Las actividades relacionadas a salidas educativas, fuera del ámbito escolar, estarán regidas bajo la Resolución Nº 1000, “Reglamento de Salidas Escolares”. 

Para poder refrendar el conocimiento adquirido en dichas salidas, será parte de la actividad, a modo de etapa de evaluación, la confección de un informe, en el cual se detallarán las actividades realizadas y las conclusiones obtenidas de la experiencia. La estructura del informe será la siguiente: 

Los trabajos prácticos de investigación para ser presentados, deben estar confeccionados en computadora según el siguiente formato: 

\begin{enumerate}[label=\alph*)]






\item Portada: datos de los alumnos participantes, docente, materia, tema, curso, año. 

\item Índice (en caso de ser necesario) 

\item Contenido – Desarrollo. (con rótulo correspondiente) 

\item Conclusiones. 

\item Ubicación, datos y características de la empresa o entidad visitada, así como el nombre de la persona que haya sido el responsable de coordinar la visita. 
 
\end{enumerate}
 

\section{Acuerdos sobre las notas y evaluaciones.} 

\subsection{ Sobre la evaluación continua}

La etapa de evaluación constituye una parte del proceso educativo. Razón por la cual, la evaluación, debe propiciar la adquisición de nuevos conocimientos y será usada con la mayor atención posible. 

Se utilizarán instancias de evaluación individual y grupal. Para tal fin se aplicarán rúbricas y evaluaciones permanentes con el objetivo de lograr un seguimiento continuo y permanente de los avances que realicen los alumnos. 

La evaluación continua les brinda a los estudiantes una serie de beneficios dentro del entorno educativo, ya que les da una mayor oportunidad de aprobar las asignaturas, teniendo en cuenta que los contenidos podrán ser asimilados y aprendidos de una manera progresiva y profunda, percibiendo todo el apoyo y colaboración de manera continua por parte del docente. 

Las posibles metodologías de evaluación continua a aplicar podrán ser: 
\begin{itemize}

\item Lecciones 

\item Evaluaciones escritas 

\item Evaluaciones virtuales

\item Evaluaciones de opción múltiple 

\item Coevaluaciones 

\item Autoevaluaciones 

\item Ejercicios de verificación 

\item Preguntas orales

\item Cuestionarios escritos 
\end{itemize}

\subsection{Sobre la evaluación integradora:}
La misma será consensuada con los alumnos, en función de la trayectoria realizada y a los tiempos disponibles. Las modalidades a implementar serán:  
\begin{itemize}
    
\item Proyectos de investigación 

\item Monografías 

\item Exposiciones orales 

\item Evaluaciones escritas 

\item Trabajos prácticos 

 
\end{itemize}

\section{Acuerdos sobre la participación de los padres. }

El padre o tutor que lo requiera, puede dialogar con el docente en el horario correspondiente al dictado de la materia. Es realmente importante la participación de los padres en el proceso de aprendizaje. Cualquier duda que necesiten evacuar podrán realizarla al docente, sin ningún tipo de problema, durante los horarios habituales de cursado. 

 

\section{Acuerdos sobre faltas e inasistencias.}

La evaluación continua como herramienta, requiere de compromiso real en las partes intervinientes, ya que la responsabilidad de evaluar en forma permanente el trabajo del alumno, tiene como contrapartida, el deber del alumno en el cumplimiento de la asistencia.  

Si por algún motivo el alumno falta el día de la evaluación escrita, avisada con anticipación, deberá comunicar al profesor antes del inicio del módulo correspondiente (por intermedio de un compañero, familiar, etc.), la imposibilidad de concurrir y cuando se reintegre presentar el certificado correspondiente, de esta forma el alumno tendrá el derecho a realizar la evaluación correspondiente cuando se reintegre a clase. El incumplimiento de este procedimiento implica la caducidad del derecho del alumno, a tener una nueva instancia de evaluación. 

 

\section{Consideraciones finales}

El presente compendio de normas tiene por objetivo que el desarrollo de la materia se realice de la mejor forma posible, dando las posibilidades a todos los alumnos de ser los partícipes indispensables en su propia formación y en la construcción de su propio conocimiento. Por lo tanto, se espera del alumno una participación positiva de cada uno y aportes significativos que le permitan alcanzar las capacidades esperadas.  

El uso de las TIC y de todos los entornos virtuales, son nuevas instancias de aprendizaje y herramientas que posibilitan un intercambio más dinámico y potenciar experiencias que enriquecerán las trayectorias de cada alumno. 
\vspace{5cm}
\begin{table}[h!]
    \centering
    \begin{tabular}{c c}
…………………………          &
…………………………   \\
FIRMA DEL ALUMNO         & FIRMA DEL PADRE O TUTOR \\
…………………………          &
…………………………   \\
ACLARACIÓN & ACLARACIÓN \\
    \end{tabular}

\end{table}


 


\end{document}